\honest{Disguised Queries}{Eliezer Yudkowsky}{2008}

You work in a factory, sorting blue eggs, called ``bleggs,'' into one bin and red cubes, called ``rubes,'' into another. Bleggs are also furred, flexible and opaque; rubes are smooth, hard and translucent. When a purple egg with fur comes along, you call it a strange blegg, though by definition a blegg is blue.

You ask Susan, the senior sorter, why the bins exist. She has no better answer than ``otherwise they'd be all mixed up,'' until she mentions that bleggs contain vanadium and rubes palladium.

So is a blegg just a ``vanadium-containing object''? No. About 2 per cent of ordinary bleggs contain palladium. Such an object goes in the rube bin, but it will glow in the dark like any blegg. Asked ``Is it a blegg?'', you should answer as if it were a rube when the question is which bin, and as if it were a blegg when the question is whether it glows. The question ``may stand in for different queries on different occasions,'' and ``If it weren't standing in for some query, you'd have no reason to care.'' \nb{This is William James's pragmatic test (1907): ``If no practical difference whatever can be traced, then the alternatives mean practically the same thing, and all dispute is idle.''}

People who argue that atheism is a religion ``because it states beliefs about God'' are really trying to argue, I think, that atheists reason no better than believers. I quote none of them. \nb{The exchange behind this is the hypothetical one in Yudkowsky's earlier post ``Uncritical Supercriticality.''} The irrational part is looking up ``atheism'' or ``religion'' in a dictionary. ``How can reality vary with the meaning of a word?'' ``The points in thingspace don't move around when we redraw a boundary.''

People often do not see that an argument about where to draw a boundary is an argument about what to infer. Hence the phrase ``disguised query.''
